\thispagestyle{empty}

\section*{Résumé}

Ce mémoire présente la conception et le déploiement d’une infrastructure Cloud privée hyperconvergée basée exclusivement sur des technologies open source. L’objectif est d’intégrer les ressources de calcul, de stockage et de réseau au sein d’une plateforme unique offrant une gestion centralisée, une haute disponibilité et une utilisation optimale des ressources. La solution proposée s’appuie sur des outils tels que Proxmox VE, Ceph et Open vSwitch afin de construire un environnement Cloud privé performant, flexible et évolutif. Ce travail met en évidence la pertinence des solutions open source pour la mise en place d’infrastructures modernes, fiables et économiques, tout en permettant une gestion unifiée et automatisée des ressources informatiques.

\vspace{0.5cm}

\section*{Abstract}

This thesis presents the design and deployment of a hyperconverged private cloud infrastructure built entirely with open-source technologies. The main objective is to combine computing, storage, and networking resources into a single platform that provides centralized management, high availability, and efficient resource utilization. The proposed solution relies on tools such as Proxmox VE, Ceph, and Open vSwitch to create a private cloud environment that is reliable, flexible, and scalable. This work demonstrates the effectiveness of open-source solutions for building modern and cost-effective infrastructures while ensuring unified and automated management of IT resources.

\vspace{0.5cm}

\section*{ملخص}

تتناول هذه المذكرة تصميم ونشر بنية تحتية لسحابة خاصة فائقة التقارب بالاعتماد الكامل على تقنيات مفتوحة المصدر. يهدف هذا العمل إلى دمج موارد الحوسبة والتخزين والشبكات ضمن منصة موحدة توفر إدارة مركزية، وتوافراً عالياً، واستغلالاً أفضل للموارد. تعتمد البنية المقترحة على أدوات مثل Proxmox VE وCeph وOpen vSwitch لبناء بيئة سحابية خاصة تتميز بالأداء والمرونة وقابلية التوسع. كما يبرز هذا العمل فعالية وموثوقية الحلول مفتوحة المصدر في إنشاء بنى تحتية حديثة واقتصادية مع ضمان إدارة موحدة وآلية لموارد تكنولوجيا المعلومات.
